\begin{vitapage}

\begin{center}
    Charles Gannon
\end{center}

\begin{edu}
  \item[2021]  \textbf{B.S.\ in Physics}, Idaho State University 
  \item[2021-present] \textbf{Graduate Student}, University of California Merced
\end{edu}

\begin{skills}
 \item \textbf{Scientific Computing} --- Simulation of physical systems, semi-analytic modeling, N-body simulations,  HPC/SLURM
 \item \textbf{ML \& Statistics} --- Regression, Classification, Pipeline Construction, PyTorch, scikit-learn, NumPy/SciPy/pandas, agentic AI 
 \item \textbf{Languages} --- Python, C++, Fortran, Bash, LaTex

\end{skills}

\begin{jobs}
 \item \textbf{Graduate Research Assistant: Computational Astrophysics} 2022 -- Present, University of California Merced
  \begin{enumerate}

      \item Designed a novel, ultra-fast hybrid semi-analytic pipeline to simulate alternate dark matter models, providing and order of magnitude speed up over traditional simulation methods
      \item Orchestrated large suites (tens of thousands of core hours) of state-of-the-art semi-analytic simulations of dark matter halos and galaxies using SLURM and HPC systems
      \item Built open source automated pipelines for analyzing and visualizing simulation output from galaxy simulations to quantify theoretical uncertainties and develop priors on dark matter properties directly relevant for observational studies
      \item Contributed to several open source scientific packages include \gal and \pyh, and authored the {\tt SubScript} python package.
      \item Communicated results through papers including published first author result, conference papers and presentations.
  \end{enumerate}
 \item \textbf{Team Lead: Data Science Challenge} Lawrence Livermore National Lab (LLNL)
  \begin{enumerate} 
      \item Lead a 3 person team for the week long challenge to build a pipeline for materials in 3D additive manufacturing 
      \item Developed an agentic pipeline for detecting and classifying defects with $> 90 \%$ accuracy
  \end{enumerate}
  \item \textbf{Teaching Assistant} University of California, Merced
  \begin{enumerate}
      \item Taught multiple (3-4) undergraduate labs and discussion sections with $20-40$ students in introductory physics classes for biologists, engineers and physicists
      \item Incorporated principles of active learning and technological aids into discussion sections and labs
  \end{enumerate}
  
\end{jobs}


\nocitepublications{2025PhRvD.112b3532G, 2025arXiv251107513G, 2025arXiv251107513G, 2024MNRAS.535.1652K, 2024PhRvD.110b3019D, 2026arXiv260405237N}
\bibliographystylepublications{unsrtnat}
\bibliographypublications{publications}

\begin{presentations}
\item[2026]%
				\textbf{American Astronomical Society, Winter 2026 Meeting} --- Contributed Talk:
				``Accelerated Core Collapse of Highly Tidally Truncated SIDM Halos''
%
\item[2024 -- 2026]%
       \textbf{UC Merced Cosmo Gal Research Symposium} --- Contributed talks: ``Dark Matter: A Lensing Perspective''
				 and
				 ``T1DES in SIDM: A New Class of Self Interacting Dark Matter Subhalo Simulations''
%
\item[2025]%
				\textbf{American Physical Society, March Meeting 2025} --- Poster:
				``Dark Matter: A Lensing Perspective''
%
\item[2024]%
				\textbf{Cosmic Signals of Dark Matter Physics: New Synergies} --- Kavli Institute of Theoretical Physics, UC Santa Barbara ---
				Contributed talk and poster: ``Dark Matter: A Lensing Perspective''
%
\item[2024]%
				\textbf{KITP Workshop} --- Kavli Institute of Theoretical Physics, UC Santa Barbara ---
				Month long workshop on Dark Matter Theory, Simulation, and Analysis in the Era of Large Surveys.
\end{presentations}

\begin{leadership}
\item\textbf{Outreach} --- Developed and taught sections of a week-long intensive for high school students covering simulation, modeling, and visualization of dark matter and gravitational lensing using Python.
\item\textbf{Awards} --- Graduated from Idaho State University with Magna Cum Laude Distinction. University of California, Merced Physics Travel Fellowship.
\end{leadership}

\end{vitapage}
